% =====================================================================
% UMT-TBS-official — BÁO CÁO CHI TIẾT CODE (LaTeX, tiếng Việt)
% File chính: biên dịch bằng pdflatex (babel vietnamese + vntex T5).
%
% Build:   make            (hoặc: pdflatex main.tex 3 lượt + cp main.pdf
%                            UMT_TBS_CodeDetail.pdf)
% Output:  UMT_TBS_CodeDetail.pdf
%
% Phạm vi: đi sâu từng folder / file / hàm / dòng theo 4 lớp IoT chuẩn
%   (Cảm biến, Mạng, Xử lý, Ứng dụng) + CI/Bảo mật + nối các lớp.
% =====================================================================
\documentclass[a4paper,12pt]{report}

% ---------- Encoding / Vietnamese ----------
\usepackage[utf8]{inputenc}
\usepackage[T5,T1]{fontenc}
\usepackage[vietnamese]{babel}

% ---------- Layout ----------
\usepackage[margin=2.5cm]{geometry}
\usepackage{graphicx}
\usepackage{booktabs}
\usepackage{longtable}
\usepackage{array}
\usepackage{amssymb}
\usepackage{xcolor}
\usepackage{hyperref}
\usepackage{float}
\usepackage{tikz}
\usetikzlibrary{shapes.geometric,arrows.meta,positioning,fit,calc}
\usepackage{fancyhdr}
\usepackage{titlesec}
\usepackage{enumitem}
\usepackage{caption}
\usepackage{listings}

\hypersetup{
    colorlinks=true,
    linkcolor=blue!50!black,
    urlcolor=blue!50!black,
    citecolor=blue!50!black,
    pdftitle={UMT-TBS-official — Báo cáo chi tiết Code (V2)},
    pdfauthor={UMT-TBS}
}

\setlength{\headheight}{15pt}
\pagestyle{fancy}
\fancyhf{}
\fancyhead[L]{Báo cáo Chi tiết Code — UMT-TBS-official (V2)}
\fancyhead[R]{\thepage}
\renewcommand{\headrulewidth}{0.4pt}

\titleformat{\section}{\large\bfseries}{\thesection.}{0.5em}{}
\titleformat{\subsection}{\normalsize\bfseries}{\thesubsection.}{0.5em}{}

\newcommand{\code}[1]{\texttt{#1}}
% Macro ảnh raster: bọc IfFileExists để build không gãy khi thiếu ảnh.
\newcommand{\pic}[2]{%
  \IfFileExists{#1}{\includegraphics[width=#2\textwidth]{#1}}{}%
}

\title{\textbf{Báo cáo Chi tiết Code\\[6pt]
\large Hệ thống Cảnh báo Va chạm Xe tải (Truck Blind-spot Warning System)\\
Phân tích theo 4 lớp IoT — folder, file, hàm, dòng code}}
\author{UMT-TBS-official}
\date{\today}

\begin{document}

\maketitle

\begin{abstract}
\noindent
Báo cáo này đi sâu vào \textbf{từng folder, từng file, từng hàm và từng dòng
code} của hệ thống \textbf{UMT-TBS-official} (Hybrid V2), tổ chức theo \textbf{4
lớp IoT chuẩn}: \emph{Cảm biến} (Perception) — 6 cảm biến JSN-SR04T và bộ lọc
trên ESP32-S3 sensor-node; \emph{Mạng} (Network) — hai đường ESP-NOW (chính,
\textless{}50\,ms) và MQTT lên CoreIoT (phụ, 100--500\,ms); \emph{Xử lý}
(Processing) — Rule-Chain ThingsBoard
(\code{cloud/coreiot/rule\_chain/supersonic\_rule\_chain.json}) và bộ tool
\code{tools/}; \emph{Ứng dụng} (Application) — màn hình LVGL v9
waveshare-screen và dashboard cloud. Mỗi chương gồm sơ đồ minh hoạ (TikZ, ảnh
trong \code{report-code/pictures/}), bảng liệt kê hàm/ký hiệu và phân tích
dòng-cuối có tham chiếu quy tắc R1--R12.
\end{abstract}

\tableofcontents
\newpage

% ================= CÁC CHƯƠNG (được viết ở các bước roadmap) =================
\IfFileExists{chapters/00_tong_quan.tex}{% =====================================================================
% Chương 0 — Tổng quan hành trình dữ liệu (4 lớp IoT chuẩn ↔ repo)
% =====================================================================

\setcounter{chapter}{-1}

\chapter{Chương 0 --- Tổng quan hành trình dữ liệu}

Chương này đặt bối cảnh cho toàn bộ báo cáo: ánh xạ \textbf{4 lớp IoT chuẩn}
(Cảm biến, Mạng, Xử lý, Ứng dụng) vào hệ thống UMT-TBS-official và vẽ sơ đồ
dòng dữ liệu trên cả hai đường truyền. Các chương sau (1--7) đi sâu vào từng
thư mục, từng file, từng hàm và từng dòng code.

\section{Bốn lớp IoT chuẩn nhìn từ hệ thống}

\begin{center}
\small
\renewcommand{\arraystretch}{1.2}
\begin{longtable}{@{}p{3.4cm}p{4.6cm}p{4.2cm}p{4.2cm}@{}}
\caption{Ánh xạ 4 lớp IoT chuẩn sang thành phần thực trong repo.}\\
\toprule
\textbf{Lớp} & \textbf{Vai trò (ẩn dụ)} & \textbf{Thành phần phần cứng/mạng} & \textbf{Thư mục chính trong repo} \\
\midrule
\endhead
\textbf{Cảm biến} (Perception) & Giác quan: đo thực tế, tạo gói dữ liệu &
6\,$\times$ JSN-SR04T + ESP32-S3 sensor-node + còi GPIO11 &
\code{firmware/shared/thresholds.h}, \code{firmware/sensor-node/src/} \\
\textbf{Mạng} (Network) & Hệ thần kinh: vận chuyển dữ liệu &
ESP-NOW (đường chính), Wi-Fi \& MQTT (đường phụ) &
\code{espnow\_protocol.h}, \code{espnow\_client.cpp}, \code{espnow\_receiver.c}, \code{plugins/coreiot} \\
\textbf{Xử lý} (Processing) & Não bộ: lưu trữ, phân loại, ra quyết định &
CoreIoT broker chạy Rule-Chain (ThingsBoard) &
\code{cloud/coreiot/rule\_chain/}, \code{tools/test\_mqtt\_coreiot.py} \\
\textbf{Ứng dụng} (Application) & Ý thức: hiển thị, tự động hoá &
Waveshare 7" RGB LVGL v9 + Dashboard cloud &
\code{firmware/waveshare-screen/src/}, \code{components/ui\_dashboard}, \code{components/coreiot\_client} \\
\bottomrule
\end{longtable}
\end{center}

\section{Hai đường truyền song song (Hybrid V2)}

Điểm khác biệt so với mô hình IoT 4 lớp "dọc" thông thường: hệ thống này có
\textbf{hai tuyến} cùng vận chuyển gói dữ liệu \code{espnow\_sensor\_msg\_t}:

\begin{center}
\begin{tikzpicture}[
  node distance=1.4cm and 2.0cm,
  box/.style={draw, rounded corners=2pt, minimum height=0.9cm, minimum width=3.2cm, font=\small, align=center},
  lbl/.style={font=\small\bfseries, align=center}
]
  % LỚP 1
  \node[lbl] (l1) at (0,0) {Lớp 1 — Cảm biến};
  \node[box, fill=blue!8] (sen) at (0,-1.1) {6$\times$ JSN-SR04T\\ESP32-S3 sensor-node\\SensorTask (100 ms)};
  % LỚP 2
  \node[lbl] (l2) at (0,-2.8) {Lớp 2 — Mạng};
  \node[box, fill=green!8] (espnow) at (-3.2,-4.0) {ESP-NOW\\kênh 1, 500 ms\\đường chính <50 ms};
  \node[box, fill=green!12] (mqtt) at (3.2,-4.0) {MQTT/Wi-Fi\\app.coreiot.io:1883\\đường phụ 100--500 ms};
  % LỚP 3
  \node[lbl] (l3) at (3.2,-5.8) {Lớp 3 — Xử lý (CoreIoT/ThingsBoard)};
  \node[box, fill=orange!12] (rc) at (3.2,-7.0) {Rule-Chain 7 node\\script JS (100/30 cm)\\Timeseries + Attributes};
  % LỚP 4
  \node[lbl] (l4) at (0,-8.8) {Lớp 4 — Ứng dụng};
  \node[box, fill=red!8] (screen) at (-3.2,-10.0) {Waveshare-screen\\LVGL v9 UI\\Cảnh báo + còi hiển thị};
  \node[box, fill=red!12] (dash) at (3.2,-10.0) {Dashboard cloud\\warning\_status\\B9 nghiệm thu};
  % ARROWS
  \draw[-{Stealth}] (sen) -- (espnow) node[midway, right, font=\small\itshape] {packed 30 B};
  \draw[-{Stealth}] (sen) -- (mqtt) node[midway, right, font=\small\itshape] {telemetry 500 ms};
  \draw[-{Stealth}] (espnow) -- (screen) node[midway, right, font=\small\itshape] {trực tiếp};
  \draw[-{Stealth}] (mqtt) -- (rc);
  \draw[-{Stealth}] (rc) -- node[midway, right, align=center, font=\small\itshape] {attributes\\notifyDevice} (screen);
  \draw[-{Stealth}] (rc) -- (dash);
\end{tikzpicture}
\captionof{figure}{Hành trình của một gói dữ liệu khoảng cách qua 2 đường truyền.}
\end{center}

\textbf{Đường chính (ESP-NOW)}: SensorTask đọc/lọc 6 cảm biến $\rightarrow$
\code{espnow\_client} gửi gói 30\,B mỗi 500\,ms $\rightarrow$
\code{espnow\_receiver} trên screen nhận, cập nhật \code{sensor\_model}
$\rightarrow$ UI đổi màu. Toàn tuyến mục tiêu dưới 50\,ms và \textbf{không phụ
thuộc Internet}.

\textbf{Đường phụ (MQTT/CoreIoT)}: \code{coreiot\_client} publish telemetry
\texttt{d1..d6 + nearest\_cm} mỗi 500\,ms $\rightarrow$ broker
\code{app.coreiot.io} $\rightarrow$ Rule-Chain chạy script JS phân loại
\code{warning\_status} $\rightarrow$ lưu Timeseries, đẩy Attributes ngược về
screen và hiển thị Dashboard.

\section{Cây repo theo lớp}

\begin{verbatim}
UMT-TBS-official/
|-- firmware/
|   |-- shared/                 [HOP DONG (R2/R3/R4) - dung cho ca 2 board]
|   |   |-- espnow_protocol.h   [Lop 2: payload ESP-NOW packed]
|   |   +-- thresholds.h        [Lop 1: cham cam bien, nguong, buzzer]
|   |-- sensor-node/            [Lop 1 + 2: đo, loc, gui]
|   |   `-- src/
|   |       |-- ultrasonic_sensor.cpp   [1: đọc Trig/Echo]
|   |       |-- distance_filter.cpp     [1: loc cluster-EMA]
|   |       |-- shared_state.cpp        [1: trang thai 6 slot]
|   |       |-- buzzer.cpp              [1: coi GPIO11]
|   |       |-- espnow_client.cpp       [2: gui ESP-NOW]
|   |       |-- main.cpp                [1: SensorTask/BuzzerTask/NetworkTask]
|   |       `-- plugins/coreiot/
|   |           `-- coreiot_client.*    [2: MQTT 500 ms]
|   `-- waveshare-screen/       [Lop 4: hien thi, nhan 2 duong]
|       |-- src/main.c + bsp/           [4: khoi dong, RGB LCD + GT911]
|       `-- components/
|           |-- espnow_receiver/        [2: nhan ESP-NOW, LINKED badge]
|           |-- sensor_model/           [4: trang thai + classify + stale]
|           |-- ui_dashboard/           [4: giao dien LVGL v9]
|           `-- coreiot_client/         [2+4: MQTT nhan attributes]
|-- cloud/coreiot/rule_chain/   [Lop 3: snapshot Rule-Chain (R11)]
|-- tools/                      [Lop 3: test MQTT, guard scripts]
|-- .github/workflows/ci.yml    [CI + bao mat]
`-- report-code/                [ban bao cao nay]
\end{verbatim}

\section{Quy ước đọc báo cáo}

\begin{itemize}
\item \code{\textbackslash code\{...\}} — tên file/ký hiệu trong code.
\item Mỗi mục hàm có bảng: \textbf{Ký hiệu} / \textbf{File} / \textbf{Lớp} /
\textbf{Vai trò}, rồi phân tích dòng-cuối có tham chiếu dòng thực tế
(\code{file.cpp:34}).
\item Gắn cờ quy tắc CONSTITUTION khi liên quan: \textbf{R2} shared contract,
\textbf{R3} thresholds một nguồn, \textbf{R4} SENSOR\_COUNT từ
\code{sizeof}, \textbf{R7} file $\le$ 400 dòng, \textbf{R11} rule-chain
snapshot.
\item Ảnh raster nếu có đặt trong \code{report-code/pictures/} (macro
\code{\textbackslash pic}); sơ đồ vector vẽ bằng TikZ ngay trong chương.
\end{itemize}}{}
\IfFileExists{chapters/01_perception.tex}{% =====================================================================
% Chương 1 — Lớp Cảm biến (Perception)
% Phân tích: shared/thresholds.h, ultrasonic_sensor.cpp, distance_filter.cpp,
%            shared_state.cpp, main.cpp (SensorTask/BuzzerTask)
% =====================================================================

\chapter{Lớp Cảm biến (Perception)}

Lớp Perception "chạm" thế giới thực: điều khiển 6 cảm biến JSN-SR04T, tạo gói
dữ liệu khoảng cách đã lọc và giữ trạng thái dùng chung cho các task. Phần mềm
nằm ở \code{firmware/sensor-node/src/} + một file hợp đồng dùng chung
\code{firmware/shared/thresholds.h}.

\section{Danh mục file và ký hiệu chính}

\begin{center}
\small
\renewcommand{\arraystretch}{1.15}
\begin{longtable}{@{}p{5.6cm}p{3.2cm}p{7.4cm}@{}}
\caption{Các file lớp Cảm biến và ký hiệu công khai.}\\
\toprule
\textbf{File} & \textbf{Lớp} & \textbf{Ký hiệu chính} \\
\midrule
\endhead
\code{firmware/shared/thresholds.h} & Chung & \code{SENSOR\_PINS}, \code{SENSOR\_COUNT}, \code{sensor\_zone\_t}, \code{BUZZER\_*}, \code{FILTER\_*} \\
\code{src/ultrasonic\_sensor.cpp} & 1 & \code{UltrasonicSensor::begin/readOnce}, ISR \code{onEchoChange} \\
\code{src/distance\_filter.cpp} & 1 & \code{process}, \code{getStable}, \code{reset} (cluster-EMA) \\
\code{src/shared\_state.cpp} & 1 & \code{sharedStateSet/Get/GetNearest} + mutex \\
\code{src/main.cpp} & 1 & \code{sensorTask} (core 1), \code{buzzerTask} (core 0) \\
\bottomrule
\end{longtable}
\end{center}

\section{Hợp đồng dùng chung: thresholds.h (124 dòng, R2/R3/R4)}

File này là \textbf{nguồn duy nhất} mọi ngưỡng/layout/chuẩn đo của cả 2 board.
Bốn khối chính:

\begin{enumerate}
\item \textbf{Zone} (dòng 32--46): \code{sensor\_zone\_t} ba mức
\code{SAFE/CAUTION/DANGER}; ngưỡng \code{SENSOR\_CAUTION\_CM 100} và
\code{SENSOR\_DANGER\_CM 30}. Cả 2 board include cùng file nên không thể lệch
(R3).
\item \textbf{Còi} (dòng 51--56): \code{BUZZER\_PIN 11} (đã chuyển khỏi GPIO48
tránh xung đột PSRAM — ghi chú dòng 96--102), phân cấp WARNING \textless{}50\,cm
(period 3000\,ms) / DANGER \textless{}20\,cm (period 1000\,ms, beep 120\,ms).
\item \textbf{Đo/lọc} (dòng 62--81): \code{SOUND\_SPEED\_CM\_PER\_US 0.0343},
\code{ECHO\_TIMEOUT\_US 40000}, \code{MEASURE\_INTERVAL\_MS 100}, bộ lọc
cluster-EMA (\code{FILTER\_HISTORY\_SIZE 9}, \code{FILTER\_MIN\_SAMPLES 5},
\code{FILTER\_EMA\_ALPHA 0.30f}, \code{FILTER\_JUMP\_CONFIRM\_COUNT 3}).
\item \textbf{Layout 6 cảm biến} (dòng 104--116): \code{static const
sensor\_pin\_cfg\_t SENSOR\_PINS[] = \{\{5,6\},...\}} — FRONT (5/6),
LEFT\_FRONT (7/8), RIGHT\_FRONT (9/10), LEFT\_REAR (17/18), RIGHT\_REAR
(21/38), REAR (3/4).
\end{enumerate}

\textbf{R4 dòng 118--123} — không gõ tay count:
\begin{verbatim}
#define SENSOR_COUNT (sizeof(SENSOR_PINS) / sizeof(SENSOR_PINS[0]))
...
TBS_STATIC_ASSERT(SENSOR_COUNT == 6, "SENSOR_COUNT must equal 6 sensor slots");
\end{verbatim}
\code{TBS\_STATIC\_ASSERT} dùng \code{static\_assert} (C++) hoặc
\code{\_Static\_assert} (C) tuy theo trình biên dịch — đó là cách 1 header
phục vụ cả 2 core (R2).

\textbf{Lưu ý ánh xạ vị trí} (ghi chú dòng 96--102): thứ tự vật lý
\code{SENSOR\_PINS} \textbf{không trùng} \code{espnow\_slot\_t} trong
\code{espnow\_protocol.h}; sensor-node giữ mảng \code{SENSOR\_ESPNOW\_SLOT[]}
trong \code{main.cpp} để đổi vật lý $\rightarrow$ wire. Báo cáo này phân tích
phía Mạng ở Chương 2.

\section{ultrasonic\_sensor.cpp (104 dòng) — đọc 1 lần đo}

\begin{center}
\begin{tikzpicture}[
  node distance=0.7cm and 1.1cm,
  box/.style={draw, rounded corners=2pt, minimum height=0.8cm, minimum width=3.6cm, font=\scriptsize, align=center},
  arr/.style={-{Stealth}, thin}
]
  \node[box, fill=blue!6] (trig) {Task phát xung 20 us\\\code{TRIGGER\_HIGH\_US}};
  \node[box, fill=blue!6, below=of trig] (isr) {ISR \code{onEchoChange} (CHANGE)\\bắt cạnh lên + cạnh xuống};
  \node[box, fill=blue!6, below=of isr] (sem) {\code{xSemaphoreGiveFromISR} $\rightarrow$ task thức dậy};
  \node[box, fill=green!6, below=of sem] (dist) {\code{durationUs * 0.0343 / 2}\\kiểm MIN/MAX (15--500 cm)};
  \node[box, fill=green!6, below=of dist] (flt) {\code{distance\_filter.process(cm)}\\cluster-EMA};
  \node[box, fill=orange!8, below=of flt] (st) {\code{sharedStateSet(i, out, valid)}\\mutex bảo vệ};
  \draw[arr] (trig) -- (isr);
  \draw[arr] (isr) -- (sem);
  \draw[arr] (sem) -- (dist);
  \draw[arr] (dist) -- (flt);
  \draw[arr] (flt) -- (st);
\end{tikzpicture}
\captionof{figure}{Chuỗi xử lý một lần đo của sensor-node.}
\end{center}

Phân tích từng hàm:

\textbf{\code{begin(trigPin, echoPin)} (dòng 19--31)}: gán chân, tạo binary
semaphore, gọi \code{attachInterruptArg(... CHANGE)} — chạy ISR trên cả hai
cạnh Echo. Truyền \code{this} làm hằng số để ISR tĩnh gọi đúng instance.

\textbf{\code{IRAM\_ATTR onEchoChange} (dòng 36--52)} — ISR, diệt nhẹ:
\begin{itemize}
\item Cạnh lên (\code{digitalRead != HIGH}): lưu \code{\_echoRiseUs = micros()}.
\item Cạnh xuống: \code{\_echoDurationUs = micros() - \_echoRiseUs}, rồi
\code{xSemaphoreGiveFromISR} + \code{portYIELD\_FROM\_ISR} nếu task cao hơn
đang chờ — đúng khuôn ISR FreeRTOS (không gọi hàm blocking/Serial).
\end{itemize}

\textbf{\code{waitEchoLow(timeoutUs)} (dòng 58--73)}: đảm bảo Echo ở LOW trước
khi Trigger mới; dùng \code{vTaskDelay(1)} thay busy-wait để nhường CPU.

\textbf{\code{readOnce()} (dòng 75--104)} — lõi đo:
\begin{enumerate}
\item \code{waitEchoLow(ECHO\_LOW\_TIMEOUT\_US)} — nếu Echo vẫn HIGH $\Rightarrow$
\code{error} trả về sớm (không đo khi Echo chưa rảnh).
\item \code{xSemaphoreTake(\_echoSemaphore, 0)} — xoá semaphore sót từ lần đo
trước.
\item Phát xung: LOW 5\,us $\rightarrow$ HIGH 20\,us $\rightarrow$ LOW.
\item \code{xSemaphoreTake(\_echoSemaphore, timeoutTicks)} — \textbf{không
busy-wait}; task ngủ đến khi ISR báo hoặc timeout 40\,ms.
\item \code{distanceCm = durationUs * SOUND\_SPEED\_CM\_PER\_US / 2.0f} —
quãng đường chia 2 (đi + về).
\item Loại vùng mù \code{< MIN\_DISTANCE\_CM (15)} và vượt tầm
\code{> MAX\_DISTANCE\_CM (500)}; lỗi trả về qua \code{reading.error} (con trỏ
hằng chuỗi — không cấp phát).
\end{enumerate}

\section{distance\_filter.cpp (368 dòng) — cluster-EMA}

Bộ lọc kết hợp \textbf{cụm} (cluster) và \textbf{EMA} để loại nhiễu siêu âm
(phản xạ phụ, vật lạ) trong \code{FILTER\_HISTORY\_SIZE = 9} mẫu gần nhất:
\begin{itemize}
\item \code{sortedMedian} — insertion sort tại chỗ (đủ nhanh với $\le 9$ phần tử,
không dùng STL theo yêu cầu dòng 6--9).
\item \code{findBestCluster} — nhóm các mẫu nằm trong dung sai động
(\code{FILTER\_BASE\_CLUSTER\_TOLERANCE\_CM 8} hoặc 8\% của tâm) và chọn cụm
đông nhất.
\item \code{getJumpThreshold / getJumpTolerance} — ngưỡng "nhảy xa" tối thiểu
30\,cm (hoặc 25\%) phải được xác nhận \code{FILTER\_JUMP\_CONFIRM\_COUNT = 3}
lần liên tiếp, tránh đổi trạng thái do nhiễu.
\item \code{process(cm)} $\rightarrow$ \code{FilterResult\{outputCm, hasOutput\}}
áp EMA \code{FILTER\_EMA\_ALPHA 0.30f}; nếu quá nhiều mẫu lỗi liên tiếp
\code{FILTER\_RESET\_AFTER\_INVALID 15} thì \code{reset()}.
\end{itemize}

Đây chính là "độ trễ 0,5--0,8\,s" được nhắc trong checklist T3.5: bộ lọc cần
$\ge$5 mẫu + 3 lần xác nhận bước nhảy — báo cáo này ghi nhận là hạn chế cần đo
thật (xem Chương 7).

\section{shared\_state.cpp (69 dòng) — trạng thái dùng chung}

Dùng mutex FreeRTOS để SensorTask/BuzzerTask/NetworkTask/CoreIoTTask chia sẻ
an toàn:
\begin{itemize}
\item \code{sharedStateInit} — tạo mutex một lần.
\item \code{sharedStateSet(i, cm, valid)} — ghi slot nếu chỉ số hợp lệ, có
mutex (timeout 10\,ms).
\item \code{sharedStateGet(i, cm)} — trả \code{true} nếu slot đang \code{valid}.
\item \code{sharedStateGetNearest(cm)} — duyệt 6 slot, trả khoảng cách \textbf{nhỏ
nhất} hợp lệ — nguồn cho BuzzerTask và CoreIoTTask (\code{nearest\_cm}).
\end{itemize}

\section{main.cpp — SensorTask và BuzzerTask}

\textbf{Khởi tạo} (\code{setup}, dòng 295--348): 4 task cố định core
(\code{xTaskCreatePinnedToCore}), điều kiện biên dịch \code{\#if USE\_COREIOT}
cho CoreIoTTask (R6 — bản coreiot dùng \code{yolo\_uno\_coreiot});
\code{loop()} gọi \code{vTaskDelete(nullptr)} vì toàn bộ logic đã ở các task.

\textbf{\code{sensorTask} (core 1, dòng 126--199)} — vòng đo 100\,ms:
\begin{itemize}
\item Đầu task: \code{begin} 6 cảm biến + \code{reset} 6 bộ lọc; chờ 500\,ms
ổn định; in mapping \code{S[0..5] Trig/Echo -> ESP-NOW slot} và cảnh báo
\code{warnIfReservedPin} nếu trùng GPIO bảo lưu (47/48 PSRAM).
\item Vòng lặp: với mỗi cảm biến $\rightarrow$ \code{readOnce()}; nếu
\code{error != nullptr} tăng \code{s\_invalidCount[i]}; đủ
\code{FILTER\_RESET\_AFTER\_INVALID (15)} lần lỗi $\rightarrow$
\code{reset() + sharedStateSet(i, 0, false)}; nếu hợp lệ
$\rightarrow$ \code{process()} rồi \code{sharedStateSet(i, output, hasOutput)}.
\item \code{vTaskDelayUntil} giữ nhịp 100\,ms \textbf{không cộng dồn trễ} và
không chặn task khác.
\end{itemize}

\textbf{\code{buzzerTask} (core 0, buzzer.cpp:63 dòng)} — đọc
\code{sharedStateGetNearest} mỗi vòng; chọn period theo ngưỡng 50/20\,cm và
toggle GPIO11 bằng \code{millis()} không blocking (xem Chương 6 máy trạng
thái còi).

\section{Kết luận lớp Cảm biến}

Đầu ra của lớp này là một mảng \code{SENSOR\_COUNT=6} giá trị \emph{đã lọc +
hợp lệ} trong \code{shared\_state}, sẵn sàng cho lớp Mạng đóng gói. Mọi ngưỡng
được định nghĩa một lần duy nhất (R3), số lượng cảm biến suy từ mảng (R4), file
giới hạn \code{ultrasonic\_sensor.cpp} 104 dòng / \code{shared\_state.cpp} 69
dòng đều $\le$ 400 (R7).}{}
\IfFileExists{chapters/02_network.tex}{% =====================================================================
% Chương 2 — Lớp Mạng (Network)
% Phân tích: espnow_protocol.h, espnow_client.cpp, networkTask (main.cpp),
%            espnow_receiver.c, coreiot_client (2 board), platformio.ini
% =====================================================================

\chapter{Lớp Mạng (Network)}

Lớp Mạng vận chuyển dữ liệu đã lọc từ Lớp 1 tới đích bằng \textbf{hai tuyến
song song}: ESP-NOW (đường chính trên xe) và MQTT/Wi-Fi (đường phụ lên cloud).
Cả hai tuyến đều đọc cùng một nguồn \code{shared\_state} nhưng gói tin và tần
suất khác nhau.

\section{Tổng quan hai tuyến}

\begin{center}
\small
\renewcommand{\arraystretch}{1.15}
\begin{longtable}{@{}p{4.0cm}p{5.6cm}p{6.0cm}@{}}
\caption{So sánh đường chính và đường phụ.}\\
\toprule
\textbf{Tiêu chí} & \textbf{ESP-NOW (chính)} & \textbf{MQTT (phụ)} \\
\midrule
\endhead
Giao thức & Link Layer Espressif, không AP & TCP/MQTT qua Wi-Fi \\
Tần suất & \code{ESPNOW\_SEND\_INTERVAL\_MS 500} & \code{COREIOT\_PUBLISH\_INTERVAL\_MS 500} \\
Payload & \code{espnow\_sensor\_msg\_t} packed 30\,B & JSON \code{d1..d6 + nearest\_cm} \\
Độ trễ mục tiêu & <50\,ms & 100--500\,ms \\
Phụ thuộc Internet & Không & Có \\
Đích & Waveshare-screen & app.coreiot.io:1883 \\
\bottomrule
\end{longtable}
\end{center}

\section{Hợp đồng dây: espnow\_protocol.h (R2)}

File duy nhất định nghĩa giao diện ESP-NOW; cấm định nghĩa lại symbol ở nơi
khác (R2 — \code{grep} toàn firmware = 1). Các hằng số quan trọng:

\begin{verbatim}
#define ESPNOW_CHANNEL 1
static const uint8_t ESPNOW_PEER_MAC[6] = {0x64,0xe8,0x33,0x7c,0x3f,0xe0};
#define ESPNOW_SEND_INTERVAL_MS 500
#define ESPNOW_LINK_TIMEOUT_MS 1500
typedef enum { ESPNOW_SLOT_FRONT=0, ESPNOW_SLOT_REAR=1,
               ESPNOW_SLOT_LEFT_FRONT=2, ESPNOW_SLOT_LEFT_REAR=3,
               ESPNOW_SLOT_RIGHT_FRONT=4, ESPNOW_SLOT_RIGHT_REAR=5 } espnow_slot_t;
typedef struct __attribute__((packed)) {
    float distance_cm[ESPNOW_SENSOR_SLOT_COUNT];
    uint8_t valid[ESPNOW_SENSOR_SLOT_COUNT];
} espnow_sensor_msg_t;
TBS_STATIC_ASSERT(ESPNOW_SENSOR_SLOT_COUNT == SENSOR_COUNT, ...);
\end{verbatim}

Bản đồ bộ nhớ gói tin 30\,B:

\begin{center}
\small
\renewcommand{\arraystretch}{1.15}
\begin{longtable}{@{}lcc@{}}
\caption{Bố cục payload ESP-NOW.}\\
\toprule
\textbf{Trường} & \textbf{Offset} & \textbf{Kích thước} \\
\midrule
\endhead
\code{distance\_cm[0..5]} & 0 & 6\,$\times$ 4\,B = 24\,B (float) \\
\code{valid[0..5]} & 24 & 6\,$\times$ 1\,B = 6\,B \\
\textbf{Tổng} & 0 & \textbf{30\,B} \\
\bottomrule
\end{longtable}
\end{center}

Lưu ý ánh xạ \textbf{vật lý $\rightarrow$ dây} nằm ở \code{main.cpp:53}:
\code{SENSOR\_ESPNOW\_SLOT[SENSOR\_COUNT]} — \code{SENSOR\_PINS[0]} (FRONT) là
dây slot FRONT, \code{SENSOR\_PINS[5]} (REAR) là dây slot REAR. Nhờ
\code{static\_assert} (R4), số slot trên dây luôn bằng \code{SENSOR\_COUNT}.

\section{espnow\_client.cpp (42 dòng) — phát ESP-NOW}

\begin{itemize}
\item \code{begin()}: \code{WiFi.mode(WIFI\_STA)} + \code{disconnect()};
\code{esp\_wifi\_set\_channel(ESPNOW\_CHANNEL,...)} cố định kênh 1; init
ESP-NOW; đăng ký \code{onDataSent}; thêm peer với MAC đích
(\code{peerInfo.encrypt = false} — không mã hoá, chấp nhận vì mạng cục bộ và
payload không nhạy cảm).
\item \code{sendReading(msg)}: \code{esp\_now\_send(ESPNOW\_PEER\_MAC, ...,
sizeof(msg))} — gửi đúng 30\,B, trả \code{true} nếu đưa vào hàng đợi thành
công.
\end{itemize}

\section{networkTask trong main.cpp — điều phối phát 500\,ms}

Task core 0 (dòng 202--240):
\begin{enumerate}
\item \code{s\_espNowClient.begin()} một lần.
\item Mỗi 500\,ms (\code{now - lastSendMs >= ESPNOW\_SEND\_INTERVAL\_MS}):
tạo \code{espnow\_sensor\_msg\_t msg = \{\}}; với mỗi cảm biến
\code{i = 0..SENSOR\_COUNT-1} đọc \code{sharedStateGet}, nếu hợp lệ thì ghi
\code{msg.distance\_cm[slot]} và \code{msg.valid[slot] = 1} theo
\code{SENSOR\_ESPNOW\_SLOT[i]}; rồi \code{sendReading}.
\item \code{vTaskDelay(20)} giữa các lần kiểm tra — không gửi blocking.
\end{enumerate}

Điểm yếu cố ý: slot không lắp phần cứng giữ \code{valid=0} $\rightarrow$
waveshare hiển thị ``--'' thay vì báo DANGER giả.

\section{espnow\_receiver.c (112 dòng) — nhận và theo dõi liên kết}

\begin{itemize}
\item \code{on\_data\_recv}: \textbf{từ chối gói sai kích thước}
(\code{len != sizeof(espnow\_sensor\_msg\_t)}) — ngăn buffer-overread; copy gói
rồi cập nhật timestamp \code{s\_slot\_rx\_us[i]} cho từng slot có
\code{valid[i]}; gọi callback driver với RSSI.
\item \code{espnow\_receiver\_init(cb)}: đặt kênh, init ESP-NOW (chấp nhận trả
\code{ESP\_ERR\_ESPNOW\_EXIST}), cài callback.
\item \code{espnow\_receiver\_is\_linked()}: so \code{now - s\_last\_any\_rx\_us}
với \code{ESPNOW\_LINK\_TIMEOUT\_MS * 1000} (1500\,ms) — nguồn cho badge
``ESP-NOW: LINKED/NO LINK'' trên UI (xem Chương 5).
\end{itemize}

\section{MQTT — CoreiotClient (sensor-node, Arduino + PubSubClient)}

File \code{plugins/coreiot/coreiot\_client.cpp} — đường phụ chỉ build khi
\code{USE\_COREIOT=1} (env \code{yolo\_uno\_coreiot}, R6):
\begin{itemize}
\item \code{begin()}: \code{setServer(COREIOT\_BROKER, COREIOT\_PORT)};
\code{WiFi.begin(WIFI\_SSID, WIFI\_PASSWORD)} — credential từ
\code{credentials.h} \textbf{sinh tự động, gitignored} (R1).
\item \code{ensureWifiConnected()}: không gọi lại \code{WiFi.begin()} nền
(chỉ tái kết nối nền) để tránh reset trạng thái đang thử.
\item \code{ensureMqttConnected()}: thử lại mỗi
\code{COREIOT\_MQTT\_RETRY\_INTERVAL\_MS = 3000}; kết nối bằng
\code{s\_mqttClient.connect("sensor-node", SENSOR\_NODE\_DEVICE\_TOKEN,
nullptr)} — \textbf{access token là username MQTT}, password rỗng (chuẩn
ThingsBoard/CoreIoT).
\item \code{publishTelemetry(payload)}: publish lên
\code{COREIOT\_TELEMETRY\_TOPIC} (\code{v1/devices/me/telemetry}).
\end{itemize}

\subsection{coreiotTask — payload telemetry 500\,ms (main.cpp:252--290)}

\begin{verbatim}
{"d1":25.0,"d2":120.0,...,"d6":400.0,"nearest_cm":25.0,"has_nearest":true}
\end{verbatim}
Dùng \code{snprintf} JSON đơn giản (không thư viện JSON Arduino), đọc từng slot
qua \code{sharedStateGetValue} + gần nhất qua \code{sharedStateGetNearest};
publish thất bại thì bỏ qua (loop() tự duy trì kết nối) — \textbf{không làm hỏng
route cảnh báo chính}.

\section{MQTT — coreiot\_client (waveshare-screen, ESP-IDF esp-mqtt)}

File \code{components/coreiot\_client/coreiot\_client.c} — client đầy đủ hơn
vì screen vừa gửi vừa nhận:
\begin{itemize}
\item \code{mqtt\_event\_handler} — khi \code{MQTT\_EVENT\_CONNECTED}:
  \begin{itemize}
  \item Subscribe \code{v1/devices/me/telemetry},
\code{v1/devices/me/attributes} và \code{v1/devices/me/attributes/response/+}
(QoS 1) — nhận cảnh báo real-time từ Rule-Chain (đường phụ).
  \item Publish telemetry \code{reboot\_event} (\{status ONLINE, reason,
wifi\_ssid\}) để dashboard biết màn hình vừa khởi động.
  \end{itemize}
\item \code{MQTT\_EVENT\_DATA}: ghi \code{s\_last\_mqtt\_rx\_time\_ms} (dùng cho
badge MQTT trên UI) và gọi \code{s\_data\_cb} cho lớp Ứng dụng.
\item \code{MQTT\_EVENT\_DISCONNECTED}: báo \code{s\_mqtt\_cb(false)}; kết nối
lại do esp-mqtt tự thử.
\end{itemize}

\section{Tập trung: luồng mạng 1 chu kỳ}

\begin{verbatim}
SensorTask (100 ms) -> shared_state (mutex)
   |-- networkTask (500 ms): khoi tao msg 30B -> esp_now_send -> screen
   `-- coreiotTask (500 ms): snprintf JSON d1..d6+nearest
              -> PubSubClient.publish(v1/devices/me/telemetry)
              -> broker app.coreiot.io:1883 -> Lớp 3 (Chương 3)
\end{verbatim}

\section{Kết luận lớp Mạng}

Mạng đóng gói đúng hợp đồng (R2), gửi đều nhịp, không blocking, không phụ
thuộc nhau: mất MQTT không ảnh hưởng ESP-NOW. Trạng thái liên kết (link timeout
1500\,ms) và thời gian MQTT rx được chuyển lên Lớp Ứng dụng (Chương 5) để hiển
thị badge.}{}
\IfFileExists{chapters/03_processing.tex}{% =====================================================================
% Chương 03 — Lớp Xử lý (Processing)
% Phân tích: tools/test_mqtt_coreiot.py, supersonic_rule_chain.json,
%            tools/guard/check_rulechain_thresholds.py
% =====================================================================

\chapter{Lớp Xử lý (Processing)}
\label{ch:processing}

Lớp Xử lý nằm giữa Lớp Cảm biến (nơi dữ liệu thô được lọc) và Lớp Ứng dụng
(nơi hiển thị cảnh báo).  Trên \textbf{sensor-node}, bản sao knn-filtered
\code{sharedStateGetNearest()} được truyền thẳng sang ESP-NOW.  Trên
\textbf{đường phụ MQTT/CoreIoT}, dữ liệu \code{d1..d6 + nearest\_cm} đi lên
broker, nơi \textbf{rule-chain} xử lý thành \code{warning\_status} rồi đẩy xuống
\code{waveshare-screen}.

\section{Tool kiểm thử MQTT (R1/R3/R11)}

\begin{center}
\begin{longtable}{@{}p{4.2cm}p{3.8cm}p{6.0cm}@{}}
  \toprule
  \textbf{Thuộc tính} & \textbf{Giá trị} & \textbf{Ghi chú} \\
  \midrule
  \endhead
  File & \code{tools/test\_mqtt\_coreiot.py} & 217 dòng, Python 3.12+ \\
  Broker & \code{app.coreiot.io:1883} & ThingsBoard CE (CoreIoT) \\
  Topic & \code{v1/devices/me/telemetry} & \code{username = access\_token} \\
  Payload keys & \code{d1..d6, nearest\_cm,} & mirror firmware/SIM/I \\
  & \code{has\_nearest, warning\_status} & (thresholds.h $\rightarrow$ classify) \\
  Token nguồn & \code{config/keys.json} & R1: \textbf{không} hardcode \\
  Dry-run & \code{--dry-run} & in payload, \textbf{không} kết nối MQTT \\
  \bottomrule
\end{longtable}
\end{center}

\textbf{Hàm chính:}
\begin{itemize}
\item \code{classify(distance)} (\code{test\_mqtt\_coreiot.py:55}): trả
  \code{"DANGER"} nếu $\leq 30$, \code{"CAUTION"} nếu $\leq 100$, ngược lại
  \code{"NORMAL"}.
\item \code{build\_payload(distance, seq)} (\code{test\_mqtt\_coreiot.py:64}):
  tạo dict với \code{d1..d6 = distance}, \code{nearest\_cm = distance},
  \code{has\_nearest = True}, \code{warning\_status = classify(...)},
  \code{vehicle\_detected}, \code{timestamp} (epoch ms), \code{seq} (tuỳ chọn).
\item \code{resolve\_token(args, cfg)} (\code{test\_mqtt\_coreiot.py:78}):
  ưu tiên \code{--token} $>$ env \code{COREIOT\_TOKEN} $>$
  \code{config/keys.json.SENSOR\_NODE\_DEVICE\_TOKEN}.
\item \code{run\_publisher(args)} (\code{test\_mqtt\_coreiot.py:133}): kết nối
  MQTT (\code{paho.mqtt}), gửi \code{json.dumps(payload)}, hỗ trợ
  \code{--loop --interval}.
\end{itemize}

\section{Sơ đồ Rule-Chain (7 nút)}
\label{sec:rulechain-graph}

\begin{center}
\begin{tikzpicture}[
  node distance=1.4cm and 1.6cm,
  box/.style={draw, rounded corners, font=\scriptsize, align=center,
               text width=3.6cm, fill=white},
  arrow/.style={-stealth, thick, gray!80},
]
% Nút 0-1: đỉnh trái
\node[box] (saveTS0)  {Save Timeseries\\(sensor-node)};
\node[box, below of=saveTS0] (saveA1) {Save Attr\\(sensor-node)};
% Nút 2: trung tâm trái
\node[box, right of=saveTS0, xshift=2.2cm, yshift=-0.7cm]
  (sw2) {Type Switch\\(MESSAGE)};
% Nút 3: trung tâm phải
\node[box, right of=sw2, xshift=2.6cm] (js3)
  {JS Transform\\(classify $\leq 30$/$\leq 100$)};
% Nút 4: phải
\node[box, right of=js3, xshift=2.8cm] (ori4)
  {Change Originator\\$\rightarrow$ waveshare-screen};
% Nút 5-6: đỉnh phải
\node[box, right of=ori4, xshift=2.8cm, yshift=0.7cm]
  (shared5) {Shared Attr\\(notifyDevice=\textbf{true})};
\node[box, below of=shared5] (saveTS6)
  {Save Timeseries\\(waveshare-screen)};

% Liên kết
\draw[arrow] (sw2) -- node[above,font=\tiny]{TS}(saveTS0);
\draw[arrow] (sw2) -- node[below,font=\tiny]{Attr}(saveA1);
\draw[arrow] (sw2) -- (js3);
\draw[arrow] (js3) -- (ori4);
\draw[arrow] (ori4) -- node[above,font=\tiny]{Shared}(shared5);
\draw[arrow] (ori4) -- node[below,font=\tiny]{TS}(saveTS6);
\end{tikzpicture}
\end{center}

\textbf{Luồng dữ liệu:} \code{sensor-node} gửi telemetry $\rightarrow$
\textbf{TypeSwitch} (\#2) $\rightarrow$ ba nhánh: (1) lưu TS sensor-node (\#0),
(2) lưu Attr sensor-node (\#1), (3) \textbf{JS Transform} (\#3) $\rightarrow$
\textbf{ChangeOriginator} (\#4) $\rightarrow$ hai nhánh: shared Attr với
\textbf{notifyDevice=true} (\#5) để \code{waveshare-screen} nhận qua MQTT, lưu
TS waveshare-screen (\#6).

\section{JS Transform — Phân loại khoảng cách}

Script JavaScript trong nút \#3 (\code{supersonic\_rule\_chain.json:79--81})
thực hiện phân loại \textbf{giống hệt} firmware:
\begin{enumerate}
\item Lấy \code{nearest\_cm} nếu \code{has\_nearest == true};
  nếu không, tìm min trong \code{[d1..d6]}.
\item Phân loại: \code{distance $\leq$ 30.0} $\rightarrow$ \textbf{DANGER};
  $\leq 100.0$ $\rightarrow$ \textbf{CAUTION}; else \textbf{NORMAL}.
\item Tạo \code{warning\_status}, \code{vehicle\_detected}, \code{relay},
  \code{buzzer}, \code{nearest\_cm}, \code{source\_device}.
\item Trả về \code{msgType: "POST\_ATTRIBUTES\_REQUEST"} — quy ước
  ThingsBoard ghi \textbf{shared attributes}, không phải timeseries.
\end{enumerate}

\noindent\textbf{Giá trị hằng số trong JS (R3):} \code{30.0} và \code{100.0}
hardcoded trong script, khớp \code{thresholds.h} \code{DANGER\_CM = 30} /
\code{CAUTION\_CM = 100}.  Script \code{check\_rulechain\_thresholds.py}
kiểm tra sự khớp này tại mỗi lần build (xem §\ref{sec:threshold-check}).

\section{Validate Thresholds (R3/R11)}
\label{sec:threshold-check}

\begin{center}
\begin{longtable}{@{}p{4.5cm}p{3.5cm}p{5.5cm}@{}}
  \toprule
  \textbf{Tool} & \textbf{Thực thi} & \textbf{Cơ chế} \\
  \midrule
  \endhead
  \code{check\_rulechain\_thresholds.py} & \code{CI + local} & Tải
  \code{thresholds.h} (grep \code{CAUTION\_CM}, \code{DANGER\_CM}) +
  \code{rule\_chain.json} (grep \code{30.0}, \code{100.0}) \\
  Exit code & \code{0 = khớp} & Nếu bất kỳ ngưỡng nào lệch: exit 1 + in
  giá trị lệch \\
  CI gate & \code{ci.yml} & Chạy \code{python3 tools/guard/
  check\_rulechain\_thresholds.py}; fail $\rightarrow$ CI đỏ \\
  \bottomrule
\end{longtable}
\end{center}

\textbf{Trích yếu kiểm tra (R3):} mọi ngưỡng zone \textbf{chỉ có một nguồn}
tại \code{firmware/shared/thresholds.h}.  Rule-chain JSON là
\textbf{snapshot} có thể đối chiếu — khi đổi ngưỡng, phải sửa
\code{thresholds.h} trước, rule-chain được import lại trên ThingsBoard console
(Cổng B0b), và \code{check\_rulechain\_thresholds.py} xác nhận khớp.

\section{Tóm tắt luồng xử lý}
\label{sec:processing-summary}

\begin{center}
\begin{tabular}{@{}lll@{}}
  \toprule
  \textbf{Điểm bắt đầu} & \textbf{Điểm kết thúc} & \textbf{Kênh} \\
  \midrule
  \code{sharedStateGetNearest()} &
    \code{espnow\_send(...)} & ESP-NOW (chính) \\
  \code{coreiotTask:} \code{json\_dump\_d1..d6} &
    Rule-chain $\rightarrow$ screen & MQTT (phụ) \\
  Rule-chain JS Transform &
    \code{Shared Attr: warning\_status} &
    MQTT notifyDevice \\
  \bottomrule
\end{tabular}
\end{center}
}{}
\IfFileExists{chapters/04_application_core.tex}{% =====================================================================
% Chương 04 — Lớp Ứng dụng (Core)
% Phân tích: waveshare-screen main.c, sensor_model.{c,h}, bsp/*rgb_lcd_port.*
% =====================================================================

\chapter{Lớp Ứng dụng — Core}
\label{ch:application-core}

Dữ liệu sau khi qua Lớp Xử lý ($\S$\ref{ch:processing}) được
\textbf{waveshare-screen} tiếp nhận và duy trì ở một \textbf{sensor model}
(container có mutex), sau đó màn hình LVGL đọc ra để vẽ.  Lớp này gồm ba phần:
\textbf{entry point} (\code{main.c}), \textbf{sensor model} (trạng thái cảm biến),
và \textbf{BSP} (khởi tạo panel LCD RGB + backlight).

\section{Entry point — \code{main.c} (221 dòng)}
\label{sec:main-c}

\textbf{Luồng khởi động} (\code{app\_main}, \code{main.c:161}):
\begin{enumerate}
\item Khởi tạo LCD RGB: \code{waveshare\_esp32\_s3\_rgb\_lcd\_init(tear\_mode,
  rotation, \&panel, \&touch)} + \code{waveshare\_rgb\_lcd\_backlight\_on()}
  (\code{main.c:169--174}).
\item Khởi tạo LVGL adapter: \code{esp\_lv\_adapter\_init(...)} với
  \code{task\_stack\_size = 12\,KB}, \code{stack\_in\_psram = true}, rồi
  \code{register\_display} (RGB, \code{use\_psram = true}) và touch
  (\code{main.c:176--196}).
\item \code{ui\_dashboard\_init()} + \code{lv\_timer\_create(espnow\_link\_watchdog\_cb,
  500, NULL)} — watchdog link (xem $\S$\ref{sec:link-watchdog})
  (\code{main.c:201--207}).
\item \textbf{Hai đường dữ liệu}, đăng ký callback:
  \begin{itemize}
  \item \textbf{Đường phụ} CoreIoT: \code{coreiot\_client\_set\_callbacks(
    on\_wifi\_status, on\_mqtt\_status, on\_coreiot\_data)} + \code{init}
    (\code{main.c:213--214}).
  \item \textbf{Đường chính} ESP-NOW: \code{espnow\_receiver\_init(on\_espnow\_rx)}
    — phải gọi \textbf{sau} \code{esp\_wifi\_start()} (đã được
    \code{coreiot\_client\_init()} thực hiện) (\code{main.c:218--221}).
  \end{itemize}
\end{enumerate}

\textbf{Callback dữ liệu:}
\begin{itemize}
\item \code{on\_espnow\_rx} (\code{main.c:117}): chạy trên WiFi task; log RSSI;
  lấy LVGL lock (timeout 100\,ms) rồi với mỗi slot \code{msg->valid[i]}:
  \code{ui\_dashboard\_update\_sensor(i, distance)}; nếu \code{valid=0}
  $\rightarrow$ \code{ui\_dashboard\_clear\_sensor(i)} (không giữ số cũ trên
  màn hình, \code{main.c:127--135}).
\item \code{on\_coreiot\_data} (\code{main.c:44}): parse JSON từ rule-chain; cập
  nhật \code{distances[6]}, \code{relay}, \code{buzzer}, \code{crossing\_hazard}.
\end{itemize}

\noindent\textbf{Chọn task an toàn:} cả hai callback đều không chạy trên LVGL
task; do đó mọi truy cập UI phải qua \code{esp\_lv\_adapter\_lock(...)} với
timeout 100\,ms (\code{main.c:35}) để không block luồng MQTT/Wi-Fi mãi mãi.

\section{Sensor Model — \code{sensor\_model.{c,h}}}
\label{sec:sensor-model}

Một container \textbf{thread-safe} đơn giản dùng mutex
(\code{SemaphoreHandle\_t s\_mutex}).

\textbf{Structure} (\code{sensor\_model.h:32}):
\texttt{sensor\_reading\_t \{ uint16\_t distance\_cm; int16\_t offset\_deg;
bool is\_stale; \}}.

\textbf{Số cảm biến không gõ tay (R2/R4):}
\begin{itemize}
\item \code{sensor\_id\_t = espnow\_slot\_t} — ánh xạ thẳng enum wire slot
  (\code{sensor\_model.h:26}).
\item \code{SENSOR\_MODEL\_COUNT = ESPNOW\_SENSOR\_SLOT\_COUNT = SENSOR\_COUNT}
  (nhờ \code{static\_assert} shared). Vị trí góc \code{k\_offsets\_deg[*]} khớp
  từng slot (FRONT=0,\, REAR=180,\, LEFT\_FRONT/REAR=-90,\,
  RIGHT\_FRONT/REAR=90, \code{sensor\_model.c:13--20}).
\end{itemize}

\textbf{API} (\code{sensor\_model.c}):
\begin{itemize}
\item \code{sensor\_model\_init} (row 22): tạo mutex, khởi tạo \code{distance=0},
  \code{is\_stale=true}.
\item \code{sensor\_model\_set\_distance(id, cm)} (row 44): clamp về
  \code{[SENSOR\_RANGE\_MIN\_CM, SENSOR\_RANGE\_MAX\_CM]} rồi ghi + set
  \code{is\_stale=false}.
\item \code{sensor\_model\_clear(id)} (row 66): \code{distance=0},
  \code{is\_stale=true} — dùng khi slot \code{valid=0}.
\item \code{sensor\_model\_get / get\_all} (row 79 / 93): bản sao an toàn.
\item \code{sensor\_model\_classify(distance)} (row 108): zone theo ngưỡng chung
  \code{SENSOR\_DANGER\_CM} / \code{SENSOR\_CAUTION\_CM} (R3) — DANGER nếu
  $< 30$, CAUTION nếu $\le 100$, SAFE nếu vượt.
\end{itemize}

\section{BSP — \code{waveshare\_rgb\_lcd\_port.c} (195 dòng)}
\label{sec:bsp}

Lớp BSP (Board Support Package) cô lập việc khởi tạo phần cứng màn hình:

\textbf{Hàm chính:}
\begin{itemize}
\item \code{waveshare\_esp32\_s3\_rgb\_lcd\_init(tear\_mode, rotation,
  \&panel, \&touch)}: cấu hình \code{esp\_lcd\_panel} cho RGB 800$\times$480
  qua esp\_lcd, thiết lập touch (nếu có), trả handle để
  \code{esp\_lv\_adapter\_register\_display} dùng.
\item \code{waveshare\_rgb\_lcd\_backlight\_on()}: bật đèn nền qua GPIO
  (PWM), cần gọi sau khi panel được init.
\end{itemize}

\noindent\textbf{Vai trò trong kiến trúc:} nhờ tách BSP, \code{main.c} không
phụ thuộc trực tiếp vào GPIO/xung nhịp của bảng Waveshare — nếu đổi board, chỉ
cần thay BSP, code ứng dụng/LVGL giữ nguyên.

\section{Tóm tắt Core}
\label{sec:core-summary}

\begin{center}
\begin{longtable}{@{}p{4.6cm}p{3.4cm}p{4.8cm}@{}}
  \toprule
  \textbf{Module} & \textbf{Dòng} & \textbf{Trách nhiệm} \\
  \midrule
  \endhead
  \code{main.c} & 221 & entry point + 2 đường dữ liệu \\
  \code{sensor\_model.c} + .h & 118 + 73 & trạng thái cảm biến + mutex \\
  \code{waveshare\_rgb\_lcd\_port.c} + .h & 195 & khởi tạo LCD RGB + backlight \\
  \bottomrule
\end{longtable}
\end{center}
}{}
\IfFileExists{chapters/05_application_ui.tex}{% =====================================================================
% Chương 05 — Lớp Ứng dụng (UI / LVGL v9)
% Phân tích: ui_dashboard.{c,h}, ui_dashboard_layout.c, ui_dashboard_system.c,
%            ui_dashboard_private.h, ui_dashboard_theme.h
% =====================================================================

\chapter{Lớp Ứng dụng — UI (LVGL v9)}
\label{ch:application-ui}

Lớp UI của \textbf{waveshare-screen} là giao diện LVGL v9 hiển thị 6 vùng
cảm biến, đánh giá nguy cơ tổng thể và trang System.  Component được tách theo
trách nhiệm (\textbf{R7} — mỗi file $\le$ 400 dòng):

\begin{center}
\begin{longtable}{@{}p{5.2cm}p{2.4cm}p{5.6cm}@{}}
  \toprule
  \textbf{File} & \textbf{Dòng} & \textbf{Trách nhiệm} \\
  \midrule
  \endhead
  \code{ui\_dashboard.c} & 331 & API công khai + \code{evaluate\_hazard} + mute \\
  \code{ui\_dashboard\_layout.c} & 335 & builder widget (header, sidebar, canvas) \\
  \code{ui\_dashboard\_system.c} & 162 & trang System + thông tin heap/uptime/token \\
  \code{ui\_dashboard\_private.h} & 75 & khai báo biến global widget \\
  \code{ui\_dashboard\_theme.h} & 42 & bảng màu + ánh xạ zone $\rightarrow$ màu \\
  \code{include/ui\_dashboard.h} & 64 & API public (hàm gọi từ \code{main.c}) \\
  \bottomrule
\end{longtable}
\end{center}

\section{Đánh giá nguy cơ — \code{evaluate\_hazard}}
\label{sec:eval-hazard}

\textbf{Trung tâm của UI} là hàm \code{evaluate\_hazard()} (\code{ui\_dashboard.c:134}):
\begin{enumerate}
\item \code{sensor\_model\_get\_all(readings)} — snapshot an toàn.
\item \textbf{Skip các cảm biến chưa có dữ liệu} (\code{readings[i].is\_stale},
  \code{ui\_dashboard.c:145}): nếu không skip, \code{distance\_cm = 0} mặc định
  sẽ phân loại \textbf{DANGER} và ghim vĩnh viễn banner OVERALL (chú thích
  \code{ui\_dashboard.c:141--144}).
\item Với cảm biến còn lại, \code{sensor\_model\_classify} theo zone chung R3;
  \code{worst = max(zone)}.
\item Cập nhật \code{s\_lbl\_hazard\_overall}: \textbf{OVERALL: DANGER /
  CAUTION / SAFE}; nếu DANGER và \code{s\_alarm\_muted} thì ghi thêm
  \code{(MUTED)} (\code{ui\_dashboard.c:153--160}).
\end{enumerate}

\noindent\textbf{Mute:} \code{mute\_btn\_cb} (\code{ui\_dashboard.c:69}) đảo
\code{s\_alarm\_muted}, cập nhật nút ("Mute Alarm"/"Unmute Alarm") và gọi lại
\code{evaluate\_hazard()} cùng \code{arc\_set\_zone}.  Nút chỉ làm \textbf{tắt
blink/animation} trên zone DANGER, banner vẫn hiện DANGER — không thể nhầm bị
mute với SAFE (\code{ui\_dashboard.c:156--158}).

\section{Bố cục widget — \code{ui\_dashboard\_layout.c}}
\label{sec:layout}

\textbf{Các builder:}
\begin{itemize}
\item \code{build\_header} (row 90): thanh đầu với title, \code{ESP-NOW} /
  \code{Wi-Fi} / \code{MQTT} status.
\item \code{build\_left\_sidebar} (row 146): các nút hành động; \textbf{lưu ý}
  nút \textbf{"Calibrate"} (\code{ui\_dashboard\_layout.c:177--181}) được tạo
  nhưng \textbf{không} có \code{lv\_obj\_add\_event\_cb} — hiện chỉ là
  placeholder trực quan (không gọi callback).
\item \code{build\_center\_canvas} (row 190): bố cục hình chiếc xe — hood,
  cabin, trunk + 6 cung \code{arc} cho mỗi vùng cảm biến
  (\code{arc\_set\_zone} / \code{arc\_set\_nodata}).
\item \code{build\_right\_sidebar} (row 253): \code{OVERALL} hazard, crossing
  risk, relay, buzzer state.
\end{itemize}

\noindent\textbf{Bảng màu zone} (\code{ui\_dashboard\_theme.h:15--23}):
\texttt{COLOR\_DANGER 0xFF1744}, \texttt{COLOR\_CAUTION 0xFFD600},
\texttt{COLOR\_SAFE 0x00C853}, \texttt{COLOR\_NODATA 0x64748B},
\texttt{COLOR\_ACCENT 0x38BDF8}, nền \texttt{COLOR\_BG 0x0F172A};
ánh xạ zone $\rightarrow$ màu ở \code{ui\_dashboard\_theme.h:39--41}.

\section{Trang System — \code{ui\_dashboard\_system.c}}
\label{sec:system}

Tab System hiển thị trạng thái thiết bị, cập nhật mỗi 2\,s qua
\code{sys\_info\_timer\_cb} (\code{main.c:129} + \code{ui\_dashboard\_system.c:87}):
\begin{itemize}
\item \textbf{Heap}: \code{esp\_get\_free\_heap\_size} /
  \code{esp\_get\_minimum\_free\_heap\_size} (KB);
\item \textbf{Uptime}: \code{format\_uptime} từ \code{esp\_timer\_get\_time}
  (h:mm:ss);
\item \textbf{FW / IDF / build / flash}.
\end{itemize}

\noindent\textbf{Bảo mật token (R1):} \code{mask\_secret} (\code{ui\_dashboard\_system.c:53})
che chuỗi token; dòng \code{ui\_dashboard\_system.c:128} chỉ hiển thị
\code{coreiot\_token\_display()} đã được che — \textbf{không} bao giờ in token
đầy đủ lên màn hình.

\section{ESP-NOW status (LINKED)}
\label{sec:linked}

Hàm \code{ui\_dashboard\_set\_espnow\_status(linked)} (\code{ui\_dashboard.c:330})
đặt nhãn \code{"ESP-NOW: LINKED"} hoặc \code{"NO LINK"}.  Được gọi từ:
\begin{itemize}
\item \code{on\_espnow\_rx} — bất kỳ frame hợp lệ $\rightarrow$ LINKED;
\item \code{espnow\_link\_watchdog\_cb} — nếu hết link (xem $\S$\ref{sec:link-watchdog})
  $\rightarrow$ NO LINK + clear các slot quá \code{ESPNOW\_LINK\_TIMEOUT\_MS}.
\end{itemize}

\section{Tóm tắt UI}
\label{sec:ui-summary}

\textbf{Chuỗi cập nhật:} \code{on\_espnow\_rx/on\_coreiot\_data}
$\rightarrow$ \code{ui\_dashboard\_update\_sensor / clear\_sensor}
($\S$\ref{sec:main-c}) $\rightarrow$ \code{arc\_set\_zone}
$\rightarrow$ \code{evaluate\_hazard()} $\rightarrow$ banner OVERALL.
Mọi vòng lặp đều nằm trên LVGL task (qua \code{esp\_lv\_adapter\_lock()}),
đảm bảo thread-safe với MQTT/Wi-Fi.
}{}
\IfFileExists{chapters/06_ci_bao_mat.tex}{% =====================================================================
% Chương 06 — CI & Bảo mật
% Phân tích: .github/workflows/ci.yml, tools/guard/{scan_secrets,
%            gen_credentials, check_size, check_rulechain_thresholds}.py,
%            .gitleaks.toml
% =====================================================================

\chapter{CI \& Bảo mật}
\label{ch:ci-security}

Hệ thống bảo vệ chất lượng và bí mật bằng \textbf{CI GitHub Actions} hai job
(\code{.github/workflows/ci.yml}) + bộ tool guard trong \code{tools/guard/}.
Chương này mô tả cách CI bắt buộc các rule R1--R12 của CONSTITUTION.

\section{Hai Job trong CI}
\label{sec:ci-jobs}

\begin{center}
\begin{longtable}{@{}p{2.8cm}p{6.0cm}p{4.4cm}@{}}
  \toprule
  \textbf{Job} & \textbf{Steps} & \textbf{Rule} \\
  \midrule
  \endhead
  \code{firmware} & Build \code{yolo\_uno} + \code{yolo\_uno\_coreiot}
    (sensor-node, R6), host test \code{native}, build waveshare-screen, size
    gate & R6, R10 \\
  \code{security} & Gitleaks, scan\_secrets, gen\_credentials --check, pytest
    tools, check\_rulechain\_thresholds, R8 assert & R1, R3, R10, R11 \\
  \bottomrule
\end{longtable}
\end{center}

\section{Job firmware — build đo được (R6/R10)}
\label{sec:ci-firmware}

\begin{enumerate}
\item \textbf{Pin PlatformIO} \code{platformio==6.1.19} và Python 3.12
  (\code{ci.yml:28--33}) — đảm bảo tái lập build (R10), tránh "về máy tôi là
  pass" (chú thích \code{ci.yml:28}).
\item \textbf{Sinh credential dummy} (\code{ci.yml:37--54}): viết
  \code{config/keys.json} với token giả \code{ci-dummy-*}, rồi chạy
  \code{gen\_credentials.py --out .../credentials.h} cho cả 2 firmware.  Nhờ
  thế env \code{yolo\_uno\_coreiot} build được khi KHÔNG có token thật — token
  thật không bao giờ vào CI (R1).
\item \textbf{Build 2 env sensor-node} (\code{yolo\_uno} + \code{yolo\_uno\_coreiot},
  \code{ci.yml:56--62}) — bảo đảm module CoreIoT được build (R6: không dead
  code).
\item \textbf{Host unit test} \code{pio test -e native}
  (\code{ci.yml:64--66}): test host DistanceFilter/thresholds (10/10).
\item \textbf{Size gate} (\code{ci.yml:74--79}) qua \code{check\_size.py}:
  \begin{itemize}
  \item sensor-node \code{firmware.bin} $\le$ 1.{2}\,MB (phân vùng app
    ~1.3\,MB);
  \item waveshare-screen $\le$ 3.7\,MB (partition factory 4\,MB $\to$ 92.5\%).
  \end{itemize}
  Chặn regression kích thước \textbf{trước} khi flash (R10, "build pass" không
  đủ).
\end{enumerate}

\section{Job security — chống leak (R1/R3/R11)}
\label{sec:ci-security-detail}

\begin{enumerate}
\item \textbf{Gitleaks 8.24.0} (\code{ci.yml:96--101}): quét full-history với
  \code{.gitleaks.toml}; exit 1 khi phát hiện secret — chặn mọi token/Wi-Fi pass
  trượt vào git.
\item \textbf{scan\_secrets.py} (\code{ci.yml:103}): quét file tracked theo
  pattern; exit 0 = sạch.
\item \textbf{gen\_credentials.py --check} (\code{ci.yml:106--122}): xác nhận
  \code{config/keys.json} hợp lệ để sinh \code{credentials.h}.
\item \textbf{pytest tools} (\code{ci.yml:124--127}): chạy
  \code{tools/guard/test\_guard.py} (16 tests) — unit test cho chính các guard.
\item \textbf{check\_rulechain\_thresholds.py} (\code{ci.yml:129--130}): đối
  chiếu ngưỡng trong rule-chain (\code{100.0}/\code{30.0}) với
  \code{thresholds.h} (R3/R11, best-effort).
\item \textbf{R8 assert} (\code{ci.yml:132--139}): \code{git ls-files} phải
  KHÔNG chứa \code{sdkconfig} hoặc \code{keys.json} — bảo vệ vĩnh viễn việc
  track nhầm credential/config build.
\end{enumerate}

\section{Bộ tool guard}
\label{sec:guards}

\begin{center}
\begin{longtable}{@{}p{5.4cm}p{2.2cm}p{5.6cm}@{}}
  \toprule
  \textbf{Tool} & \textbf{Dòng} & \textbf{Vai trò} \\
  \midrule
  \endhead
  \code{scan\_secrets.py} & 122 & quét pattern secret trong file tracked \\
  \code{gen\_credentials.py} & 85 & sinh \code{credentials.h} từ keys.json \\
  \code{check\_size.py} & 72 & size-gate firmware.bin \\
  \code{check\_rulechain\_thresholds.py} & 78 & đối chiếu ngưỡng R3/R11 \\
  \code{test\_guard.py} & 214 & 16 unit test cho các guard \\
  \bottomrule
\end{longtable}
\end{center}

\noindent\textbf{Luồng bí mật (R1):}
\code{config/keys.json} (gitignored, token do chủ tài khoản tự sinh)
$\rightarrow$ \code{gen\_credentials.py} $\rightarrow$
\code{credentials.h} (gitignored) $\rightarrow$ firmware.  Literal token/
password bị \textbf{cấm} trong mọi file tracked — cổng Gitleaks + scan\_secrets
ở CI là tuyến phòng thủ cuối.
}{}
\IfFileExists{chapters/07_ket_noi.tex}{% =====================================================================
% Chương 07 — Kết nối end-to-end
% Phân tích: truy vết "55 cm" qua 2 đường; ma trận file ↔ lớp;
%            chuỗi xác minh (CI, guard).
% =====================================================================

\chapter{Kết nối end-to-end}
\label{ch:endtoend}

Chương cuối truy vết một phép đo cụ thể xuyên suốt 4 lớp và tổng kết ma trận
file$\leftrightarrow$lớp để người đọc (sinh viên/hội đồng) định vị nhanh từng phần của hệ thống.

\section{Truy vết một mẫu đo: 55\,cm}
\label{sec:trace}

Xét sensor FRONT (GPIO 5/6, \code{SENSOR\_PINS[0]}) đo được \textbf{55\,cm} — thuộc
zone CAUTION ($100 \ge 55 > 30$, R3).

\textbf{Đường chính — ESP-NOW:}
\begin{enumerate}
\item ISR \code{IRAM\_ATTR onEchoChange} bắt pulse echo
  ($\S$\ref{ch:perception}, \code{sensor\_reading}) $\rightarrow$
  \code{distance = pulse * 0.0343 / 2} $\approx$ 55.3\,cm.
\item \code{DistanceFilter} làm sạch (median + cluster-EMA) rồi
  \code{sharedStateSetNearest(55, ...)}.
\item \code{networkTask} (mỗi 500\,ms) đọc \code{sharedStateGetNearest}, ghi
  \code{msg.distance\_cm[SLOT\_FRONT]=55} + \code{valid=1} theo
  \code{SENSOR\_ESPNOW\_SLOT[0]} ($\S$\ref{ch:network}) $\rightarrow$
  \code{espnow\_send(...)} 30\,B.
\item \code{waveshare-screen} nhận frame (\code{espnow\_receiver}), gọi
  \code{on\_espnow\_rx}: \code{ui\_dashboard\_update\_sensor(FRONT, 55)}
  $\rightarrow$ \code{arc\_set\_zone(CAUTION)} + \code{evaluate\_hazard()}
  $\rightarrow$ banner \textbf{OVERALL: CAUTION} (nếu là vùng xa nhất)
  ($\S$\ref{ch:application-ui}).
\end{enumerate}

\textbf{Đường phụ — MQTT/CoreIoT:}
\begin{enumerate}
\item \code{coreiotTask} (500\,ms) đóng \code{\{d1:55.3,...,has\_nearest:true\}}
  publish lên \code{v1/devices/me/telemetry}.
\item Rule-chain: TypeSwitch $\rightarrow$ JS Transform ($\S$\ref{ch:processing}):
  \code{55 $\le$ 100} $\rightarrow$ \code{warning\_status = "CAUTION"},
  \code{vehicle\_detected = true} $\rightarrow$ ChangeOriginator $\rightarrow$
  Shared Attr \code{notifyDevice=true}.
\item \code{waveshare-screen} đăng ký nhận attributes; parse JSON
  \code{on\_coreiot\_data} ($\S$\ref{sec:main-c}) $\rightarrow$ cập nhật
  sensor + trạng thái relay/buzzer lên UI.
\end{enumerate}

\noindent\textbf{Hai đường đều quy về cùng một \code{warning\_status} ngưỡng
duy nhất} (\code{thresholds.h}, R3) — giảm lệch dữ liệu giữa hiển thị ESP-NOW
(low-latency) và cloud.

\section{Ma trận file $\leftrightarrow$ lớp}
\label{sec:matrix}

\begin{center}
\begin{longtable}{@{}p{5.4cm}p{2.0cm}p{5.8cm}@{}}
  \toprule
  \textbf{File đường chính} & \textbf{Lớp} & \textbf{Vai trò} \\
  \midrule
  \endhead
  \code{ultrasonic\_sensor.cpp} & Cảm biến & ISR + readOnce \\
  \code{distance\_filter.cpp} & Cảm biến & lọc nhiễu, cluster-EMA \\
  \code{espnow\_client.cpp} & Mạng & gửi ESP-NOW \\
  \code{espnow\_receiver.c} & Mạng & nhận ESP-NOW (screen) \\
  \code{sensor\_model.c} & Xử lý /Core & trạng thái cảm biến + mutex \\
  \code{ui\_dashboard.c} & UI & evaluate\_hazard + mute \\
  \code{ui\_dashboard\_layout.c} & UI & bố cục widget \\
  \bottomrule
\end{longtable}

\begin{longtable}{@{}p{5.4cm}p{2.0cm}p{5.8cm}@{}}
  \toprule
  \textbf{File đường phụ} & \textbf{Lớp} & \textbf{Vai trò} \\
  \midrule
  \endhead
  \code{plugins/coreiot/coreiot\_client.} & Mạng & MQTT sensor-node \\
  \code{components/coreiot\_client/} & Mạng & MQTT screen (ESP-IDF) \\
  \code{tools/test\_mqtt\_coreiot.py} & Xử lý & gửi telemetry thử \\
  \code{supersonic\_rule\_chain.json} & Xử lý & 7-nút rule-chain \\
  \code{ci.yml} + \code{tools/guard/} & CI/Bảo mật & build + guard \\
  \bottomrule
\end{longtable}
\end{center}

\section{Chuỗi xác minh}
\label{sec:verify-chain}

\begin{itemize}
\item \textbf{Build:} \code{pio run -e yolo\_uno} / \code{yolo\_uno\_coreiot}
  + \code{pio run -e yolo\_uno} (waveshare-screen) — CI job \code{firmware}.
\item \textbf{Host test:} \code{pio test -e native} — DistanceFilter 10/10;
  \code{pytest tools/guard/test\_guard.py} — 16/16.
\item \textbf{Bảo mật:} Gitleaks (\code{.gitleaks.toml}) + \code{scan\_secrets.py}
  exit 0; \code{gen\_credentials.py --check}.
\item \textbf{Ngưỡng:} \code{check\_rulechain\_thresholds.py} — CI log
  \code{CHECK-RULECHAIN OK}.
\item \textbf{Kích thước file:} mọi \code{*.c/*.cpp/*.h} (firmware + chương báo
  cáo) $\le$ 400 dòng (R7).
\end{itemize}

\noindent Kết quả CI run gần nhất xanh trên cả 2 job; báo cáo này đối chiếu
code thật tại commit (xem \code{docs/logs/CODE\_DETAIL\_REPORT\_LOG.md}).
}{}

\end{document}